\documentclass[10pt,a4paper]{amsart}
\usepackage{amsmath,amssymb,amsthm,graphicx}
% For commutative diagrams we use Xy-pic
\usepackage[all]{xy}

% We use 2cell for 2-commutative diagrams.
\xyoption{2cell}
% Theorem environments.
%
\theoremstyle{plain}
\newtheorem{theorem}[subsection]{Theorem}
\newtheorem{proposition}[subsection]{Proposition}
\newtheorem{lemma}[subsection]{Lemma}

\theoremstyle{definition}
\newtheorem{definition}[subsection]{Definition}
\newtheorem{example}[subsection]{Example}
\newtheorem{exercise}[subsection]{Exercise}
\newtheorem{situation}[subsection]{Situation}

\theoremstyle{remark}
\newtheorem{remark}[subsection]{Remark}
\newtheorem{remarks}[subsection]{Remarks}

\numberwithin{equation}{subsection}

% Macros
%
\def\lim{\mathop{\mathrm{lim}}\nolimits}
\def\colim{\mathop{\mathrm{colim}}\nolimits}
\def\Spec{\mathop{\mathrm{Spec}}}
\def\Hom{\mathop{\mathrm{Hom}}\nolimits}
\def\Ext{\mathop{\mathrm{Ext}}\nolimits}
\def\SheafHom{\mathop{\mathcal{H}\!\mathit{om}}\nolimits}
\def\SheafExt{\mathop{\mathcal{E}\!\mathit{xt}}\nolimits}
\def\Sch{\mathit{Sch}}
\def\Mor{\mathop{\mathrm{Mor}}\nolimits}
\def\Ob{\mathop{\mathrm{Ob}}\nolimits}
\def\Sh{\mathop{\mathit{Sh}}\nolimits}
\def\NL{\mathop{N\!L}\nolimits}
\def\CH{\mathop{\mathrm{CH}}\nolimits}
\def\proetale{{pro\text{-}\acute{e}tale}}
\def\etale{{\acute{e}tale}}
\def\QCoh{\mathit{QCoh}}
\def\Ker{\mathop{\mathrm{Ker}}}
\def\Im{\mathop{\mathrm{Im}}}
\def\Coker{\mathop{\mathrm{Coker}}}
\def\Coim{\mathop{\mathrm{Coim}}}

% Boxtimes
%
\DeclareMathSymbol{\boxtimes}{\mathbin}{AMSa}{"02}

%
% Macros for moduli stacks/spaces
%
\def\QCohstack{\mathcal{QC}\!\mathit{oh}}
\def\Cohstack{\mathcal{C}\!\mathit{oh}}
\def\Spacesstack{\mathcal{S}\!\mathit{paces}}
\def\Quotfunctor{\mathrm{Quot}}
\def\Hilbfunctor{\mathrm{Hilb}}
\def\Curvesstack{\mathcal{C}\!\mathit{urves}}
\def\Polarizedstack{\mathcal{P}\!\mathit{olarized}}
\def\Complexesstack{\mathcal{C}\!\mathit{omplexes}}
% \Pic is the operator that assigns to X its picard group, usage \Pic(X)
% \Picardstack_{X/B} denotes the Picard stack of X over B
% \Picardfunctor_{X/B} denotes the Picard functor of X over B
\def\Pic{\mathop{\mathrm{Pic}}\nolimits}
\def\Picardstack{\mathcal{P}\!\mathit{ic}}
\def\Picardfunctor{\mathrm{Pic}}
\def\Deformationcategory{\mathcal{D}\!\mathit{ef}}
\usepackage{cleveref}
\crefname{originalreference}{original reference}{original references}
\Crefname{originalreference}{Original reference}{Original references}
\makeatletter
\newlabel{lemma-ses-associated-to-open}{{095L}{source}{}{sourcepointer-1}{}}
\newlabel{lemma-ses-associated-to-open@cref}{{[originalreference][0][]095L}{[1][0][]source}{}{}{}}
\makeatother
\crefname{originalreference}{original reference}{original references}
\Crefname{originalreference}{Original reference}{Original references}
\makeatletter
\newlabel{proposition-finite-higher-direct-image-zero}{{03QP}{source}{}{sourcepointer-2}{}}
\newlabel{proposition-finite-higher-direct-image-zero@cref}{{[originalreference][0][]03QP}{[1][0][]source}{}{}{}}
\makeatother
\crefname{originallemma}{original lemma}{original lemmas}
\Crefname{originallemma}{Original lemma}{Original lemmas}
\makeatletter
\newlabel{sites-cohomology-lemma-apply-Leray}{{0733}{source}{}{sourcepointer-3}{}}
\newlabel{sites-cohomology-lemma-apply-Leray@cref}{{[originallemma][0][]0733}{[1][0][]source}{}{}{}}
\makeatother
\title{Finiteness of the cohomology of a constructible sheaf on a curve is unchanged by replacing it by its extension by zero from a dense open}
\author{The Stacks Project Authors}
\date{stacks-project:etale-cohomology.tex:local-snapshot-20261011; GNU FDL 1.2 or later}
\begin{document}
\maketitle
\noindent\textit{Editorial note.} All delivered mathematical text is the author's own, copied from the pinned Stacks Project chapter file etale-cohomology.tex; no statement or proof step was rewritten, reordered or removed.\par\medskip
\noindent\textit{Editorial note.} The counted proof is the author's complete argument: the identification of the complement of the dense open as a finite set of rational points, the computation of the cohomology of the sheaf supported on that finite set through the vanishing of the higher direct images and the Leray spectral sequence, and the final long exact sequence showing the equivalence of the finiteness statements.\par\medskip
\noindent\textit{Editorial note.} The label lemma-restrict-to-open-coefficients names the counted statement, and the proof environment immediately following it in the source is delivered verbatim as its proof.\par\medskip
\noindent\textit{Editorial note.} The counted statement is carried by the author's own lemma environment in the chapter source.\par\medskip
\noindent\textit{Editorial note.} Pre-open gate: python operations/agent-authoring/tools/precard\_gate.py --source etale-cohomology.tex --label lemma-restrict-to-open-coefficients --start 15926 --end 15970 --local-id 3563 returned RESULT=OK (L3563 free, occupancy FREE, public interval clear). Candidate taken from lane-s4/halfB-priority.json.\par\medskip
\noindent\textit{Editorial note.} The statement and proof are copied from the pinned chapter file; the cited inputs are kept as external pointers with their canonical Stacks tags.\par\medskip
\noindent\textit{Allowed standard prerequisites (editorial).}\begin{itemize}
\item The short exact sequence associated to a dense open. Same chapter, the lemma giving the short exact sequence 0 -> j\_!j\textasciicircum{}\{-1\}F -> F -> i\_*i\textasciicircum{}\{-1\}F -> 0 for the inclusion of the complement of a dense open. Supplies the short exact sequence whose long exact cohomology sequence is analysed.
\item Vanishing of the higher direct images of a finite morphism. Same chapter, the proposition that a finite morphism has vanishing higher direct images. Supplies the identification of the cohomology of i\_*i\textasciicircum{}\{-1\}F over X with the cohomology of i\textasciicircum{}*F over the finite set Z via the Leray spectral sequence.
\item The Leray spectral sequence for a morphism of sites. Cohomology on Sites, the lemma applying the Leray spectral sequence to compute cohomology over the source as cohomology over the target of the pushforward. Supplies the equality H\textasciicircum{}q(X, i\_*i\textasciicircum{}\{-1\}F) = H\textasciicircum{}q(Z, i\textasciicircum{}*F).
\end{itemize}\medskip
\ifdefined\hypertarget\hypertarget{sourcepointer-1}{}\fi
\noindent\textit{Source reference.} Original reference 095L: sources/stacks-etale-cohomology/pinned-original-source.tex, source lines 11588--11588. Prerequisite pointer invoked at proof line(s) 15943; it names a standard fact of the same chapter that the counted argument uses, and every step of the counted proof built on it is retained verbatim in the delivered body.\par\medskip
\ifdefined\hypertarget\hypertarget{sourcepointer-2}{}\fi
\noindent\textit{Source reference.} Original reference 03QP: sources/stacks-etale-cohomology/pinned-original-source.tex, source lines 8245--8245. Prerequisite pointer invoked at proof line(s) 15947; it names a standard fact of the same chapter that the counted argument uses, and every step of the counted proof built on it is retained verbatim in the delivered body.\par\medskip
\ifdefined\hypertarget\hypertarget{sourcepointer-3}{}\fi
\noindent\textit{Source reference.} Original lemma 0733: sources/stacks-chow/tags.txt, source lines 9047--9047. Prerequisite pointer invoked at proof line(s) 15949; it names a standard fact of a sibling chapter that the counted argument uses, and every step of the counted proof built on it is retained verbatim in the delivered body.\par\medskip
\begin{lemma}
\label{lemma-restrict-to-open-coefficients}
Let $\Lambda$ be a Noetherian ring, let $k$ be an algebraically closed field,
let $X$ be a separated finite type scheme over $k$ of dimension $\leq 1$,
let $\mathcal{F}$ be a constructible sheaf of $\Lambda$-modules on $X_\etale$,
and let $j : X' \to X$ be the inclusion of a dense open subscheme. Then
$H^q_\etale(X, \mathcal{F})$ is a finite $\Lambda$-module if and only if
$H^q_\etale(X, j_!j^{-1}\mathcal{F})$ is a finite $\Lambda$-module.
\end{lemma}

\begin{proof}
Since $X'$ is dense, we see that $Z = X \setminus X'$ has dimension $0$
and hence is a finite set $Z = \{x_1, \ldots, x_n\}$ of $k$-rational points.
Consider the short exact sequence
$$
0 \to j_!j^{-1}\mathcal{F} \to \mathcal{F} \to i_*i^{-1}\mathcal{F} \to 0
$$
of Lemma \ref{lemma-ses-associated-to-open}. Observe that
$H^q_\etale(X, i_*i^{-1}\mathcal{F}) = H^q_\etale(Z, i^*\mathcal{F})$.
Namely, $i : Z \to X$ is a closed immersion, hence finite,
hence we have the vanishing of $R^qi_*$ for $q > 0$ by
Proposition \ref{proposition-finite-higher-direct-image-zero}, and
hence the equality follows from the Leray spectral sequence
(Cohomology on Sites, Lemma \ref{sites-cohomology-lemma-apply-Leray}).
Since $Z$ is a disjoint union of spectra of algebraically closed
fields, we conclude that $H^q_\etale(Z, i^*\mathcal{F}) = 0$
for $q > 0$ and
$$
H^0_\etale(Z, i^{-1}\mathcal{F}) =
\bigoplus\nolimits_{i = 1, \ldots, n} \mathcal{F}_{x_i}
$$
which is a finite $\Lambda$-module $\mathcal{F}_{x_i}$ is finite due to the
assumption that $\mathcal{F}$ is a constructible sheaf of $\Lambda$-modules.
The long exact cohomology sequence gives an exact sequence
$$
0 \to
H^0_\etale(X, j_!j^{-1}\mathcal{F}) \to
H^0_\etale(X, \mathcal{F}) \to
H^0_\etale(Z, i^{-1}\mathcal{F}) \to
H^1_\etale(X, j_!j^{-1}\mathcal{F}) \to
H^1_\etale(X, \mathcal{F}) \to 0
$$
and isomorphisms $H^0_\etale(X, j_!j^{-1}\mathcal{F}) \to
H^0_\etale(X, \mathcal{F})$ for $q > 1$. The lemma follows easily from this.
\end{proof}

\end{document}
